\chapter{Adquisición digital y conversión de datos}
\section{\obj}
\capacidad
\begin{itemize}
\item Leer las entradas digitales del NI myDAQ y mostrarlas como un arreglo de booleanos.
\item Convertir un arreglo de booleanos en un número entero.
\item Convertir un número entero en una cadena hexadecimal.
\end{itemize}

\section{\mat}
\textbf{A suministrar por la Escuela:}
\begin{itemize}
\item 1 dispositivo de adquisición de datos NI myDAQ
\item 1 interruptor rotativo de 16 posiciones

\end{itemize}
\textbf{A suministrar por el estudiante:}
\begin{itemize}
\item 1 breadboard
\item Cables de interconexión (jumpers) macho-macho
\end{itemize}

\section{\pro}
\begin{enumerate}
\item Realice las conexiones del circuito tal y como se indica en la Figura \ref{fig:L2F1}.

\item Cree un instrumento virtual (\emph{.vi}) que realice lo siguiente:
    \begin{itemize}
        \item Leer en un While 1 muestra del valor de las entradas DIO0, DIO1, DIO2 y DIO3 cada \SI{500}{\milli\second} y mostrarla como una matriz de booleano
        \item Convertir el valor binario obtenido en un número entero y mostrarlo en el panel frontal
        \item Dividir el número entero obtenido entre el valor máximo (15) y mostrar el valor en un ``Slide''
        \item Convertir el número entero obtenido en una cadena hexadecimal
        \item Mostrar un mensaje en el panel frontal que diga ``El selector está en XX'' donde XX es la cadena hexadecimal obtenida.
        \item Guardar la cadena hexadecimal obtenida en una matriz cada vez que exista un cambio de valor
    \end{itemize}
\begin{figure}[H]
    \tikzset{demux/.style={muxdemux, muxdemux def={Lh=4, Rh=3, NL=4, NB=0, NR=1, w=2}}}
    \tikzset{switch16/.style={muxdemux, muxdemux def={Lh=3, Rh=3, NL=3, NB=0, NR=3, w=2}}}
    \centering
    \begin{circuitikz}
        \draw (0,0) node[demux] (m){\rotatebox{90}{\small multiplex}}
        ;
        \foreach \i in {0,...,3}
        {
            \pgfmathtruncatemacro{\pin}{\i + 1}
            \draw (m.blpin \pin) node[above left]{\small DI0\i};
        }
        \draw (-4,2) node[switch16,rotate=90] (s){\rotatebox{-90}{\small switch}}
        (s.brpin 3) node[above right]{\small 1}
        (s.brpin 2) node[above right]{\small C}
        (s.brpin 1) node[above right]{\small 4}
        (s.blpin 3) node[below right]{\small 8}
        (s.blpin 2) node[below right]{\small C}
        (s.blpin 1) node[below right]{\small 2}
        ;
        \draw
        (0,2) node[ground]{}
            to[V]
        (0,4) node[above]{\SI{5}{\volt}}
            --
        (-3,4)
            -|
        (s.brpin 2)
        ;
        \draw[red]
        (m.blpin 1) --++ (-2,0) --++ (0,2.5) -| (s.brpin 3)
        ;
        \draw[blue]
        (m.blpin 2) -| (s.blpin 1)
        ;
        \draw[orange]
        (m.blpin 3) --++ (-5,0) --++ (0,3.5) -| (s.brpin 1)
        ;
        \draw[brown]
        (m.blpin 4) -| (s.blpin 3)
        ;
        \draw[dashed,blue]
        (-1.6,6) -- (1,6)node[midway, below]{myDAQ} -- (1,-1.5) -- (-1.6,-1.5) -- cycle;
    \end{circuitikz}
    \caption{Lectura de estados de un interruptor rotativo de 16 posiciones.}
    \label{fig:L2F1}
\end{figure}
\item Una vez realizado el instrumento virtual, córralo, realice al menos 10 cambios y guarde los datos de la matriz de cadenas hexadecimales como un archivo de Excel.
\item Incluya en su informe lo siguiente:
    \begin{itemize}
        \item Una captura de pantalla del diagrama de bloques
        \item Una captura de pantalla del panel frontal
        \item Una tabla con los valores de las cadenas hexadecimales obtenidas.
    \end{itemize}
\end{enumerate}

Detalle del interruptor rotativo:
\begin{center}
    \includegraphics[width=0.3\textwidth]{rotary_switch.jpg}
\end{center}
